% Anisotropy AI laminate engine manual.
% Build with:  uv run --project fea python docs/manual/build.py
% The files in generated/ are written by build.py from the repository's documents, code and
% engine results. Edit those sources, not the generated files.

\documentclass[11pt,letterpaper]{article}

\usepackage{fontspec}
\usepackage{amsmath,amssymb}
\usepackage[margin=1in,headheight=14pt]{geometry}
\usepackage{xcolor}
\usepackage{microtype}
\usepackage{titlesec}
\usepackage{enumitem}
\usepackage{booktabs,array,xltabular}
\usepackage{listings}
\usepackage{fancyhdr}
\usepackage{parskip}
\usepackage{hyperref}

\definecolor{accent}{HTML}{9A5A00}
\definecolor{ink}{HTML}{1C1F22}
\definecolor{muted}{HTML}{5D636A}
\definecolor{codebg}{HTML}{F3F4F1}
\definecolor{passgreen}{HTML}{2F6B3F}

\hypersetup{colorlinks=true, linkcolor=ink, urlcolor=accent, citecolor=ink,
  pdftitle={Anisotropy AI Laminate Engine Manual}, pdfauthor={Anisotropy AI}}

\titleformat{\section}{\sffamily\Large\bfseries\color{ink}}{\color{accent}\thesection}{0.7em}{}
\titleformat{\subsection}{\sffamily\large\bfseries\color{ink}}{\color{accent}\thesubsection}{0.6em}{}
\titleformat{\subsubsection}{\sffamily\normalsize\bfseries\color{ink}}{\thesubsubsection}{0.5em}{}
\titlespacing*{\section}{0pt}{2.2ex plus 1ex}{1.2ex}
\titlespacing*{\subsection}{0pt}{2ex plus 1ex}{0.8ex}

\setlist{nosep, leftmargin=1.4em, topsep=0.3ex}
\setlength{\emergencystretch}{3em}  % lets lines with long file names break instead of overflowing
\newcolumntype{Y}{>{\raggedright\arraybackslash}X}
\renewcommand{\arraystretch}{1.25}

\lstdefinestyle{engine}{
  language=Python, basicstyle=\ttfamily\footnotesize, columns=fullflexible, keepspaces=true,
  numbers=left, numberstyle=\tiny\color{muted}, numbersep=7pt,
  backgroundcolor=\color{codebg}, frame=leftline, framerule=2pt, rulecolor=\color{accent},
  xleftmargin=2.8em, framexleftmargin=2.2em, aboveskip=0.4ex, belowskip=1.2ex,
  keywordstyle=\color{accent}\bfseries, commentstyle=\color{muted}\itshape,
  stringstyle=\color{passgreen}, showstringspaces=false, breaklines=true}
\lstnewenvironment{enginecode}[1][]{\lstset{style=engine,#1}}{}

% \codefrom{function}{file}{lines}: caption above a code listing.
\newcommand{\codefrom}[3]{\par\smallskip\noindent{\sffamily\footnotesize\color{muted}%
  Code: \textcolor{ink}{\textbf{#1}} in #2, lines #3}\par\nopagebreak}
% \adrmeta{status}{date}{url} and \adrpart{label}: decision record layout.
\newcommand{\adrmeta}[3]{{\sffamily\small\color{muted}Status: #1\quad Date: #2\quad \href{#3}{Source file}}\par}
\newcommand{\adrpart}[1]{\par\smallskip\noindent{\sffamily\footnotesize\bfseries\color{accent}\MakeUppercase{#1}}\par\nopagebreak}

\input{generated/meta}

\pagestyle{fancy}
\fancyhf{}
\renewcommand{\headrulewidth}{0pt}
\fancyhead[L]{\sffamily\footnotesize\color{muted}Anisotropy AI \textperiodcentered{} Laminate Engine Manual \textperiodcentered{} engine \EngineVersion}
\fancyhead[R]{\sffamily\footnotesize\color{muted}\thepage}
\fancyfoot[C]{\sffamily\scriptsize\color{muted}Internal use. Not for design decisions until coupon tests are compared with predictions.}

\begin{document}

\hypersetup{pageanchor=false}
\begin{titlepage}
  \sffamily
  {\color{accent}\bfseries\small ANISOTROPY AI}\par
  \vspace*{3cm}
  {\fontsize{30}{36}\selectfont\bfseries\color{ink} Laminate Engine Manual}\par
  \vspace{0.6cm}
  {\Large\color{muted} Theory, verification, decisions and limits}\par
  \vspace{2.5cm}
  {\normalsize
  \begin{tabular}{@{}ll@{}}
    Engine version & \EngineVersion \\
    Source commit & \texttt{\GitCommit} \\
    Built & \BuildDate \\
    Engine tests & \TestSummary \\
    Finite element layer & \FeaVersion, tests: \FeaTestSummary \\
  \end{tabular}}\par
  \vfill
  \colorbox{codebg}{\parbox{\dimexpr\linewidth-2\fboxsep}{\small\color{ink}\textbf{Status: internal use.}
    The engine reproduces textbook and answer-key results (Section~\ref{sec:verification}). It has not yet been compared
    with our own coupon tests, so it is not approved for design or certification decisions.}}\par
  \vspace{0.8cm}
  {\small\color{muted} This manual is generated from the repository: its prose and equations come from the
  Markdown documents, its code listings from the engine source, and every number in its tables from running the
  engine while it was built. To change it, edit those sources and rebuild (see \texttt{docs/manual/README.md}).}
\end{titlepage}
\hypersetup{pageanchor=true}

\tableofcontents
\clearpage

\section{About this manual}

This manual collects the project's documentation about the laminate engine into one readable document. Its
sections answer the questions set out in the documentation guide,
\texttt{docs/\allowbreak documentation-\allowbreak guide.md}.

\begin{xltabular}{\linewidth}{lYY}
  \toprule
  \textbf{Section} & \textbf{Question} & \textbf{Source in the repository} \\
  \midrule
  \ref{sec:built} & What have we built? & \texttt{docs/00\_system\_overview.md} \\
  \ref{sec:theory} & How is each result calculated? & \texttt{engine/docs/theory.md} and the engine code \\
  \ref{sec:verification} & Does it work, and how do we know? & \texttt{engine/validation/cases} and \texttt{docs/evaluations} \\
  \ref{sec:fea} & How are finite element models built and solved, and do they match? & \texttt{fea/docs/theory.md}, the \texttt{fea} code, \texttt{fea/validation/cases} and \texttt{docs/evaluations/fea} \\
  \ref{sec:decisions} & Why was it built this way? & \texttt{docs/decisions} \\
  \ref{sec:limits} & What could go wrong? & \texttt{docs/00\_system\_overview.md} and \texttt{docs/risks/risk\_register.md} \\
  \ref{sec:example} & What does a real calculation look like? & Computed by the engine while this manual was built \\
  \bottomrule
\end{xltabular}

Each equation in Section~\ref{sec:theory} is followed by the code that computes it, copied from the engine with its
real line numbers. Values marked as engine results were calculated during the build, so the manual cannot drift from
the code it describes.

\section{What we built}\label{sec:built}
\input{generated/overview}

\section{How results are calculated}\label{sec:theory}
\input{generated/theory}

\section{Does it work?}\label{sec:verification}
\input{generated/verification}

\section{Finite element layer}\label{sec:fea}
The finite element layer reads the Abaqus input decks the team builds for coupon models and solves them with our
own laminate membrane kernel on the engine's stiffness, so the two can be compared number for number
(decision record ADR-0029 in Section~\ref{sec:decisions}). Its equations follow, each with the code that computes
it, under the same rule as the engine's: every function the text names is listed from the source with its real line
numbers. The validation tables after the theory were produced by running the kernel while this manual was built;
agreement with the team's Abaqus runs is recorded separately once their results exist.
\input{generated/fea_theory}
\input{generated/fea_verification}

\section{Why it was built this way}\label{sec:decisions}
These are the project's decision records: the context, the options considered, the choice, its evidence, and what
would make us revisit it. The full records, including consequences, are in \texttt{docs/decisions}.
\input{generated/decisions}

\section{Limits and risks}\label{sec:limits}
\input{generated/limits}

\appendix

\section{Worked example: tension coupon}\label{sec:example}
\input{generated/example}

\section{Unit systems}\label{sec:units}
\input{generated/units}

\section{References}
\begin{itemize}[leftmargin=1.2em, itemsep=0.6ex]
  \item I. M. Daniel and O. Ishai, \emph{Engineering Mechanics of Composite Materials}, 2nd edition. Oxford University
    Press, 2006. Lamina relations, Chapter 4; lamina strength and Tsai-Wu, Chapter 6; laminate theory and engineering
    properties, Chapter 7; first-ply failure, Chapter 9.
  \item K.-J. Bathe, \emph{Finite Element Procedures}, 2nd edition. Watertown, MA: K.-J. Bathe, 2014. Isoparametric
    element formulation and numerical integration, Chapter 5 (finite element layer, Section~\ref{sec:fea}).
  \item S. P. Timoshenko and J. N. Goodier, \emph{Theory of Elasticity}, 3rd edition. McGraw-Hill, 1970. Plane stress,
    Chapter 2 (isotropic validation cases of the finite element layer).
  \item NASA-STD-7009A, \emph{Standard for Models and Simulations}. NASA, 2016.
  \item ASME V\&V 10, \emph{Standard for Verification and Validation in Computational Solid Mechanics}. ASME.
  \item NIST AI 100-1, \emph{Artificial Intelligence Risk Management Framework (AI RMF 1.0)}. NIST, 2023.
  \item Project documentation guide and its sources: \texttt{docs/documentation-guide.md}.
\end{itemize}

\end{document}
